% Generado a partir de calculus/Final_Bien/evidencia_P8.m (simulador corregido).
\section{Configuración estable}

En la configuración repulsiva actúa la bobina inferior; con $x$ medida hacia arriba, la distancia entre el imán y esa bobina es $a+x$ y la fuerza electromagnética es $u/[b(a+x)^4]$ y el modelo linealizado tiene la matriz de estado $\bigl[\begin{smallmatrix}0&1\\-4g/(a+x_1^*)&-c_1\end{smallmatrix}\bigr]$, cuyos polos son complejos con parte real negativa: $-4{,}28\pm21{,}49j$ en $x_1^*=1$~cm y $-4{,}28\pm20{,}24j$ en $x_1^*=2$~cm (factor de amortiguamiento cercano a $0{,}2$). Se evaluaron en esta configuración tres controladores diseñados durante el trabajo: el PI de la comparación de la tabla~\ref{tab:pi_pii}, una versión preliminar del PII y el PII definitivo $C(s)$ de la sección~\ref{sec:pii}. Con la planta linealizada en $x_1^*=1{,}5$~cm, los tres estabilizan el lazo, con márgenes de fase de $50{,}7^\circ$, $51{,}0^\circ$ y $43{,}3^\circ$, respectivamente. Es decir, el mismo controlador diseñado para la configuración inestable estabiliza también la estable, aunque con un margen de fase menor.

\begin{table}[ht]
\centering
\caption{Configuración estable: escalón de 1 a 2~cm con zona muerta, fricción de Karnopp y $u\in[0;\,3{,}5]$~V.}
\label{tab:p8_estable}
\begin{tabular}{|l|c|c|c|}
\hline
 & PI & PII preliminar & PII (tesis) \\ \hline
Margen de fase (lineal, $x_1^*=1{,}5$~cm) & $50{,}7^\circ$ & $51{,}0^\circ$ & $43{,}3^\circ$ \\ \hline
Sobrepaso (lineal) & $19{,}7$\,\% & $18{,}6$\,\% & $26{,}8$\,\% \\ \hline
Sobrepaso (no lineal con AD) & $10{,}5$\,\% & $11{,}1$\,\% & $9{,}8$\,\% \\ \hline
Entrada en la banda de $\pm5$\,\% & $0{,}24$~s & $0{,}36$~s & $0{,}35$~s \\ \hline
Ciclo límite pico a pico & $0{,}018$~cm & $0{,}023$~cm & $0{,}020$~cm \\ \hline
\end{tabular}
\end{table}

La figura~\ref{fig:p8_estable} y la tabla~\ref{tab:p8_estable} muestran que los tres controladores llevan el levitante de 1 a 2~cm y que, como en la configuración atractiva, la zona muerta y la fricción estática producen un ciclo límite en estado estacionario, aquí de menor amplitud. En el transitorio la señal de control alcanza brevemente la saturación de $3{,}5$~V.

\begin{figure}[ht]
    \centering
    % Estilo común de las figuras de evidencia P8 (pgfplots).
\pgfplotsset{
  pocho/.style={
    width=0.92\textwidth, height=5.2cm,
    grid=major, grid style={gray!25},
    tick label style={font=\small}, label style={font=\small}, legend style={font=\small},
    /pgf/number format/use comma,
    scaled ticks=false,
    y tick label style={/pgf/number format/fixed},
    table/col sep=comma, unbounded coords=jump,
    every axis plot/.append style={line join=round},
  },
}
\definecolor{tray}{RGB}{31,94,166}
\definecolor{refc}{RGB}{40,40,40}
\definecolor{ctl2}{RGB}{196,96,40}
\definecolor{ver}{RGB}{46,133,64}

\begin{tikzpicture}
\begin{axis}[pocho, name=principal, xmin=0, xmax=40, ymin=0.8, ymax=2.3, xlabel={},
    ylabel={posición [\si{cm}]}, legend style={at={(0.99,0.03)}, anchor=south east}, legend cell align=left]
  \addplot[ctl2, line width=0.8pt] table[x=t, y=x_pi] {images/p8_results/tikz/estable.csv}; \addlegendentry{PI}
  \addplot[tray, line width=0.8pt] table[x=t, y=x_pii] {images/p8_results/tikz/estable.csv}; \addlegendentry{PII (tesis)}
\end{axis}
\begin{axis}[pocho, at={(principal.south west)}, anchor=north west, yshift=-0.45cm, height=4.4cm,
    xmin=0, xmax=40, ymin=-0.2, ymax=3.8, xlabel={tiempo [\si{s}]}, ylabel={$u$ [\si{V}]}]
  \addplot[ctl2, line width=0.7pt] table[x=t, y=u_pi] {images/p8_results/tikz/estable.csv};
  \addplot[tray, line width=0.7pt] table[x=t, y=u_pii] {images/p8_results/tikz/estable.csv};
  \addplot[red!70!black, densely dotted, line width=0.8pt] coordinates {(0,3.5) (40,3.5)};
\end{axis}
\end{tikzpicture}

    \caption[Configuración estable]{Regulación en la configuración estable (repulsiva) con el PI y con el PII de la tesis, escalón de 1 a 2~cm. Arriba, posición; abajo, señal de control.}
    \label{fig:p8_estable}
\end{figure}

\FloatBarrier
